\documentclass[10pt,aspectratio=169]{beamer}
\usepackage{ragged2e,tikz,mathtools,amsmath,amsthm,amssymb,subcaption,pgfplots,longtable,actuarialsymbol,actuarialangle,array}
\usetheme{Unesa-Light}

\title{Judul File Disini}
\subtitle{Tulis subjudul disini. Subjudul maksimal tiga baris.}
\author{Dimas Maulana}
\date{\today}
\institute{Program Studi Sarjana Sains Aktuaria}

\begin{document}

\begin{frame}[plain]
	\titlepage
\end{frame}

\section{Introduction}
\begin{frame}
	\frametitle{A Frame}
\begin{itemize}
\item First thing
	\begin{itemize}
	\item small point
	\item fine print
	\end{itemize}
\item Second thing
	\begin{enumerate}
	\item point 1
	\end{enumerate}
\item Third thing
	\begin{description}
	\item[Research] the scientific pursuit for knowledge
	\end{description}
\end{itemize}
\end{frame}

\subsection{Text}
\begin{frame}{Another Frame}
Lorem ipsum dolor sit amet, consectetur adipisicing elit, sed do eiusmod tempor incididunt ut labore et dolore magna aliqua. Ut enim ad minim veniam, quis nostrud exercitation ullamco laboris nisi ut aliquip ex ea commodo consequat.
\end{frame}

\subsection{Blocks}
\begin{frame}[allowframebreaks]
	\frametitle{Blocks}
	
	\begin{definition}[Greetings]
		Hello World
	\end{definition}
	
	\begin{theorem}[Fermat's Last Theorem]
		$a^n + b^n = c^n, n \leq 2$
	\end{theorem}
	
	\begin{block}{What've got so far}
		\begin{enumerate}
			\item We know that ..
			\item As far as we discussed ..
		\end{enumerate}
	\end{block}
	
	\begin{alertblock}{Attention!}
		By the pricking of my thumbs.
	\end{alertblock}
	
	\begin{exampleblock}{Case 1:}
		Something wicked this way comes.
	\end{exampleblock}
	
\end{frame}

\closingframe

\end{document}